\chapter{Implementation and Validation}
\label{chapter4}

\section{Implementation}

This chapter describes the implemented system that realises the methodology defined in Chapter~\ref{chapter3}. The emphasis is on the concrete modules, scripts, and validation checks used on the deployed execution path.

\subsection{Predictor Training Pipeline}

The predictor training pipeline (\texttt{train\_predictor.py}) materialises the locked predictor configuration from Chapter~\ref{chapter3} into per-node regression artifacts from phase-labelled telemetry. In the final implementation it writes ridge-model bundles with window \(w=5\) and horizon \(h=3\).

The core implementation choices are:

\begin{itemize}
\item \textbf{Leakage-safe scaling.} A temporal split is applied before fitting the scaler, so future statistics cannot influence training features.
\item \textbf{Phase-aware sample weighting.} CPU and mixed stress targets receive higher sample weight because prediction quality during stressed operation is more important than during quiescent operation.
\item \textbf{Post-fit calibration.} Per-feature affine correction is learned on a calibration tail of the training data and applied at inference time.
\item \textbf{Artifact traceability.} Each model directory stores per-node model bundles, scalers, and a \texttt{training\_summary.json} with split sizes and error metrics.
\end{itemize}

The output of this stage is a self-contained model directory consumed later by the runtime ranker.


\subsection{Runtime Predictor}

The runtime predictor in \texttt{predictor.py} provides the inference interface consumed by the live ranker. \texttt{NodePredictor} manages a per-node sliding buffer, while \texttt{ClusterPredictor} aggregates all worker-specific predictors and exposes cluster-wide prediction.

The scalar predicted-load score is a weighted combination of CPU, RAM, and network forecasts, with CPU and memory receiving the highest weight. Two implementation details are particularly important for the final online experiment:

\begin{itemize}
\item \textbf{Warm-start fallback.} Until the sliding window is full, the ranker falls back to a bounded load estimate derived from current observations.
\item \textbf{Missing-artifact fallback.} If a node does not have a trained model bundle, the cluster predictor records it as a fallback-only node rather than failing the entire ranking path.
\end{itemize}

In the final authoritative run all 120 custom decisions used the trained-model path, but the bounded fallback path remains available for short histories or missing artifacts.


\subsection{Anomaly Monitors}

The anomaly-monitor layer in \texttt{hybrid\_scheduler.py} conforms to a common \texttt{update()}/\texttt{ready()}/\texttt{risk()} interface. The deployed runtime uses the NSA monitor selected in Chapter~\ref{chapter3}; z-score and K-means implementations remain available only for controlled offline comparison.

The NSA monitor estimates node-local \textit{self} behaviour from recent history and samples detectors in non-self regions. At inference time, risk is obtained from nearest-detector distance through a bounded logistic mapping.


\subsection{Hybrid Scheduler Engine}

The \texttt{HybridScheduler} class combines prediction and anomaly monitoring into a single ranking decision. The final live policy exposes several behaviours that matter directly to the online protocol:

\begin{itemize}
\item \textbf{Adaptive prediction/anomaly weighting.} The prediction weight remains dominant under low anomaly risk but can shift toward the anomaly term when the monitored node state deteriorates.
\item \textbf{Workload-aware capacity pressure.} Each candidate node receives an additive capacity penalty derived from the workload request and the node's allocatable CPU and memory.
\item \textbf{Observed-load guardrails.} Explicit contention and unsafe-node penalties are applied from the observed relative CPU level, and unsafe-node avoidance prevents the ranker from choosing overloaded nodes when safer alternatives exist.
\item \textbf{Decision diagnostics.} The returned score structure records base predicted load, projected load, anomaly risk, request fractions, observed relative CPU, capacity penalty, and prediction provenance.
\end{itemize}

These fields are persisted in the decision artifacts and later consumed by the online analysis pipeline.


\subsection{Offline Replay Audit}

The offline replay evaluator (\texttt{offline\_policy\_evaluation.py}) provides the audit surface for comparing prediction-only, anomaly-only, and hybrid variants on governed telemetry. The retained live family used NSA, warm-up 24, anomaly history 45, detector budget 120, detector radius 0.9, base weights 0.8 and 0.2, adaptive weighting, and explicit contention and unsafe-node penalties.

These settings were frozen before the final online run, and the final scheduler claim remains the online comparison rather than the replay audit \cite{cawley2010overfitting}.


\subsection{Authoritative Custom Decision Path}

The final online comparison attributes the custom decision directly to the implemented Python ranker. For each custom-arm pod, the runner executes the sequence summarised in Table~\ref{tab:ch4-custom-path}.

\begin{table}[h]
\centering
\small
\begin{tabular}{lp{9cm}}
\hline
\textbf{Stage} & \textbf{Implementation action} \\
\hline
Prepare candidates & Export live allocatable node capacities and enforce the common eligible worker set shared with the baseline arm. \\
Refresh live state & Seed or update the custom telemetry history with one-second samples and capture a one-shot snapshot immediately before the decision. \\
Rank and persist & Invoke \texttt{rank\_live.py} with the locked predictor, NSA configuration, workload request, and exported capacities, then persist the ranking JSON for audit. \\
Bind authoritatively & Write the chosen node into \texttt{spec.nodeName} before metric collection begins. \\
\hline
\end{tabular}
\caption{Authoritative custom decision path used in the final online comparison.}
\label{tab:ch4-custom-path}
\end{table}

Because the same worker allow-list is enforced for both arms, any observed difference is attributable to the decision rule rather than to a different feasible set.


\subsection{Online Experiment Orchestration}

The final online experiment is orchestrated by \texttt{run\_stage\_b\_authoritative\_rank\_and\_pin.sh}. Around the per-pod ranking path, the runner exports capacities at run start, derives the eligible worker pool, counterbalances arm order by trial, rotates stress profiles and stress targets across the worker set, extracts workload requests from the pod manifests, and hands enriched per-decision records to the metric logger and summary script.

As a result, the final Stage B artifacts include baseline and custom scheduling CSVs, decision snapshots, ranking JSONs, node-capacity exports, run metadata, and a summary JSON that links aggregate outcomes back to decision-level evidence.


\section{Validation and Threats to Validity}

\subsection{Predictor Validation}

Predictor validation followed the locked development/holdout protocol defined in Chapter~\ref{chapter3}. The \((w{=}5, h{=}3)\) artifact remained ahead of the naive-last and window-average baselines on the governed holdout.

The final trained artifact also produced comparable error across worker nodes, supporting the use of one common scoring surface across the live worker pool.


\subsection{Offline Replay Validation}

Offline replay validation applied the governed split to the prediction-only, anomaly-only, and hybrid candidates and fixed the safety-first hybrid family before live execution. This stage validates the locked configuration only; it does not provide the dissertation's final scheduler claim.


\subsection{Online Pipeline Readiness Validation}

Before executing the final authoritative experiment, the readiness checks in Table~\ref{tab:ch4-readiness} were completed.

\begin{table}[h]
\centering
\small
\begin{tabular}{lp{8.8cm}}
\hline
\textbf{Check} & \textbf{Implementation evidence} \\
\hline
Runner validation & Shell syntax checks passed for the Stage B orchestration scripts. \\
Python validation & \texttt{py\_compile} passed for the ranking, replay, and orchestration modules touched by the final protocol. \\
Ranking smoke check & Synthetic validation confirmed that \texttt{rank\_live.py} runs with eligible-node filtering and fallback behaviour enabled. \\
Analysis smoke check & The summary pipeline ran successfully on an existing smoke dataset. \\
Artifact readiness & The collector, metric logger, and analysis scripts accepted and preserved the per-decision fields required for telemetry audit. \\
\hline
\end{tabular}
\caption{Readiness checks completed before the final authoritative online run.}
\label{tab:ch4-readiness}
\end{table}

These checks established pipeline correctness before the final result was generated.


\subsection{Online Experiment Data Integrity}

The final authoritative run produced complete and internally consistent artifacts.

Both arms recorded 120 pods, corresponding to 10 trials with 12 pods per arm per trial. Phase coverage was complete at 30 pods per phase in both arms and 10 pods per workload within each phase. Only \texttt{k3s-worker-2}, \texttt{k3s-worker-3}, and \texttt{k3s-worker-4} appear in the final placement summaries; no Raspberry Pi placements occur in the authoritative result. The custom arm achieved a 100\% expected-node match rate, and its decision provenance shows \texttt{model: 120}, confirming that every authoritative decision was generated by the trained predictor path. Mean scheduling latency remained 0.0667 seconds in both arms, indicating that artifact capture and ranking did not introduce measurable overhead in aggregate.

Together, these checks support the interpretation that the final summary is complete and directly traceable from per-decision artifacts to aggregate results.


\subsection{Threats to Validity and Mitigations}

\textbf{Internal validity threats.} Live cluster state can drift between arms, telemetry endpoints can fail transiently, and nominally homogeneous workers can still differ through background activity. These threats were mitigated by a common eligible worker set, explicit warm-up, one-shot pre-decision snapshots, persisted decision artifacts, counterbalanced arm order, and fixed washout periods.

\textbf{Construct validity threats.} The final online experiment measures placement safety, anomaly exposure, contention, timing, fairness, and utility under controlled stress injection. These remain proxies for wider edge quality of service, but they are closer to actual placement decisions than predictor error or detector F1 alone.

\textbf{External validity threats.} The cluster remains modest relative to production edge deployments, and the final comparison is restricted to a homogeneous worker subset. This improves causal clarity for the baseline comparison but does not exhaust the heterogeneity question for mixed-capacity edge fleets.

\textbf{Selection bias threat.} Offline replay influenced the choice of the core hybrid family, but the live configuration was fixed before the authoritative run. This reduces optimistic bias from adapting the method after observing online outcomes \cite{cawley2010overfitting}.


\subsection{Reproducibility Package}

Reproducibility is supported through code, configuration, and artifact traceability, consistent with general principles for computational research \cite{peng2011reproducible}. The package comprises campaign plans and per-run telemetry artifacts, locked model directories with training summaries, offline replay audit configuration and outputs, workload manifests and metric-collection helpers, the authoritative online runner and node-capacity export, per-decision online artifacts (snapshot CSVs, ranking JSONs, metadata, and enriched arm CSVs), and the final run summary JSON used for the aggregate results in Chapter~\ref{chapter5}.

This package supports both code-level reproducibility and decision-level auditability, which is particularly important once the custom arm is evaluated through authoritative pre-decision binding.
